\documentclass[../main.tex]{subfiles}

\begin{document}

\chapter{Pemodelan Masalah Komputasi menggunakan Logika Predikat}

\begin{subcpmk}
  \item Mahasiswa mampu memodelkan masalah komputasi yang tak terstruktur di alam verbal manusia dengan merepresentasikannya menggunakan kerangka ekspresif logika predikat.
\end{subcpmk}

\section{Pemodelan Sintaks *Query* dan Struktur Data}

Kemampuan menerjemahkan kebutuhan (*requirements*) menjadi model matematika diskrit merupakan fondasi arsitektur sistem Basis Data (*Relational Databases*) modern (merujuk pada *Relational Calculus*) maupun rekayasa *Knowledge-Based AI*. Logika Predikat menyumbang sintaks formalnya secara masif.

Ketika programmer memformulasikan sintaks kueri penarikan *record*, misalnya SQL:
\lstset{language=SQL}
\begin{lstlisting}
SELECT Email FROM Pengguna WHERE Umur > 18 AND Status = 'Aktif';
\end{lstlisting}
Secara diam-diam, sang programmer tengah merancang instruksi komputasi bertumpu evaluasi model Logika Predikat.
Fungsi komputasional di \textit{backend database engine} memanifestasikan perintah tabel kolom tersebut menjadi pencarian yang memuaskan variabel kuantifikasi himpunan relasional.

Domain $U$ = {Himpunan Entitas Tabel Pengguna Aktif}.
Model Klausa Relasi Pemodelannya: $\{ x \in Pengguna \mid Umur(x) > 18 \land Status(x, Aktif) \}$.
Di sini predikat \textit{Umur} adalah fungsi pemetaan variabel tunggal iteratif ($P(x)$) dan \textit{Status} adalah relasional argumen komputatif variabel ganda $R(x, y)$.

\section{Pemodelan Logika Bersarang (*Nested Quantifiers*)}
Masalah kompleks komputasi memerlukan Kuantor jamak tak merdeka (Bersarang).
Misal pemodelan arsitektur enkripsi data komputatif: "Semua kata sandi, pasti di-enkripsi algoritma hash oleh sekurang-kurangnya sebuah server keamanan valid."

Domain $P$ = Sandi, Domain $S$ = Server. Predikat Enkripsi $E(x, y)$ (sandi x di hash server y).
Model Relasi Formal Komputasinya: $\forall x (P(x) \rightarrow \exists y (S(y) \land E(x, y)))$

\begin{contoh}
Menguraikan \textit{Deadlock Dependency}:
Misalkan $T_i$ merepresentasikan \textit{Thread/Process} CPU, dan $R_a$ tipe *Resources Data*. $W(T_i, R_a)$ artinya Thread $i$ mem-blok menunggu \textit{Resource $a$}. Masalah *System Deadlock* dimodelkan terjadi bila: "Ada thread siklis yang saling blok satu dan yang lainnya abadi." Model pendeteksi awal yang di-skrip engineer Sistem Operasi mengecek apakah *status flag komputasionalnya* merujuk \textit{True} pada relasi: $\exists x \exists y (W(x, y) \land W(y, x))$.
\end{contoh}

\begin{aktivitas}
  \item Pilih \textit{Rule} atau \textit{Guideline Form Validasi Password Sistem}. (Contoh: "Setiap sandi user setidaknya mengandung satu modifikasi atribut *special character*"). Formulasikan model kuantor berindeks logika predikatnya secara akurat di atas kanvas/kertas untuk divisualisasikan bersama sejawat Anda!
\end{aktivitas}

\begin{asesmen}
Akurasi pendefinisian ruang himpunan *Universe (U)* pada tugas perancangan spesifikasi Logika *Database Relational* yang diberikan saat ujian portofolio terstruktur.
\end{asesmen}

\begin{checklist}
  \item Mampu mengenali peranan $P(x)$ di dalam sintaks bahasa kueri \texttt{SQL} / pengkondisian data komputasi praktis *Array Filtering*.
  \item Lancar mentransformasi kebutuhan verbal rumit \textit{System User} menjadi persamaan bersarang kuantor mutlak tanpa ambiguitas terjemahan.
\end{checklist}

\end{document}
